\documentclass{itatnew}
%% !!!dolezite: ak pisete po slovensky alebo po cesky pouzite
%% \documentclass[slovensky]{itatnew}
%% \documentclass[cesky]{itatnew}
\def\OP#1{\textcolor{purple}{OP: \textit{#1}}}
\def\todo#1{\textcolor{purple}{OP: \textit{#1}}}
% \def\todo#1{\todo[noprepend,color=yellow!30]{TODO: #1}}

\begin{document}

\title{Recurrent neural networks for dialog state tracking}

\author{Petr Bělohlávek \and Vojtěch Hudeček \and
Josef Válek \and Ondřej Plátek \and TODOFilip}

\institute{Charles University in Prague,\\
\email{oplatek@ufal.mff.cuni.cz},\\ 
\texttt{http://ufal.mff.cuni.cz/ondrej-platek}}

\maketitle              % typeset the title of the contribution

\begin{abstract}
This paper discuss models for dialog state tracking using recurrent neural networks (RNN).
We present experiments on standard dialog state tracking (DST) dataset DSTC2\cite{todo}.
On one hand, RNN models became state of the art in DST,
on the other hand most state-of-the-art models are only turn-based and require preprocessing specific to evaluated dataset (e.g. DSTC2) to achieve state-of-the-art results.
We implemented three architectures which can be used in incremental settings and requires almost no preprocessing.
We compare their performance to the benchmarks on DSTC2 and discuss their properties.
With only a trivial preprocessing we were able to beat a baseline, but our models are not competitive with state-of-the-art.
\end{abstract}
%
\section{Introduction}
%
The dialog state tracking (DST) is a standard and important task for evaluating conversational agents\cite{dstc1, dstc2, dstc3, HIS model}.
The dialog state tracker summarizes hidden information state (HIS) \cite{todo} of users goal from the conversation history.
Users goals are expressed in a~formal language typically represented as a dialog act item (DAI) $(actionType, slotName, slotValue)$.
It was shown that the better dialog state tracking of HIS the conversation agents are able to achieve the better success rate they have in overall completion of the task oriented conversation.\cite{Steve Yound, Henderson,..Zilka, Alex}
The dialog state tracking translates ambiguous natural language into formal language which is convenient for reasoning and accessing external knowledges, both important aspects of successful conversation in task oriented dialog.

Most state-of-the-art system in DST have reported their performance on Dstc2 dataset\cite{dstc2henderson}. 
The full dataset is freely available since January 2014 and contains X dialogues in training set, Y dialogues in dev set and Z dialogues test set.
The conversation are annotated at turn level where the hidden information state was annotated manually in for of $(actionType, slotName, slotValue)$
according a domain ontology.
The ontology of Dstc2 captures a~restaurant domain and was also manually designed.
The dataset also contains a database of restaurants and their properties.

Our experiments use only history \todo{ASR or } or gold transcriptions of conversation history and information of database values to predict dialog state which simplify setup and still obtain reasonable results.
We argue that using only database values instead of full ontology is natural, because the system can inform users only about the values in the database.
In addition we show in Section \ref{todo} it does not harm the performance on the dataset.

In our experiments we focus only on the {\it goal} slots predictions because the other groups are trivial to predict\footnote{The slots {\it Requested} and {\it Method} have accuracies 0.95 and 0.95 on the set according to state-of-the-art\cite{JWilliams}.} and such obtaining annotated data for new domain is significantly simpler.
We show that limiting ourself to {\it goal} slots which track the properties of restaurant wanted by the user  and in particular to their values which are only in the database, we do not loose much accuracy.
See Section~\ref{todo} for details.
Tracking only constraints of the restaurants in the dialog state allow us to annotate the dialog state as query to database, which is much simpler interface for human annotators than selecting many options from Dstc2 ontology.
Simplifying the dialog state while maintaining the accuracy allows simpler portability of proposed models to new domains.

We also show that Dstc2 dataset suffers from data disbalance between training, development and test set.
The Dstc2 test set was collected in \todo{different scenario} \cite{dstc2henderson}.
Since the data of training, development and test set are distributed differently the resulting performance between training and test accuracy is rather high. 
Our experiments showed that by random resplitting of the Dstc2 data one obtain much better results.
As a result, we think that Dstc2 might suggest too pessimistic view of state-of-the-art methods in dialog state tracking because of the data distribution mismatch.

Our contribution is two fold. 
First, we compare three different architectures using RNN for dialog state tracking in Section~\ref{todo}.
Second, we analyze the Dstc2 dataset and propose changes for further data collections of similar datasets.


\section{Models}

\section{Experiments}

\subsection{Labels representation evaluation}

\begin{table}
\caption{This is the example table taken out of {\it The
\TeX{}book,} p.\,246}
\begin{center}
\begin{tabular}{r@{\quad}rl}
\hline
\multicolumn{1}{l}{\rule{0pt}{12pt}
                   Year}&\multicolumn{2}{l}{World population}\\[2pt]
\hline\rule{0pt}{12pt}
8000 B.C.  &     5,000,000& \\
  50 A.D.  &   200,000,000& \\
1650 A.D.  &   500,000,000& \\
1945 A.D.  & 2,300,000,000& \\
1980 A.D.  & 4,400,000,000& \\[2pt]
\hline
\end{tabular}
\end{center}
\end{table}


\begin{figure}
\includegraphics[width=\textwidth]{rnn}
\caption{The RNN encodes dialog history}
\end{figure}

\subsection{Data preparation experiments}

\section{Related work}

\section{Conclusion}




\paragraph{Acknowledgement}
We would like to thank nvidia, gauk, filipgrant.




%
% ---- Bibliography ----
%
\begin{thebibliography}{5}
%
\bibitem {clar:eke}
Clarke, F., Ekeland, I.:
Nonlinear oscillations and
boundary-value problems for Hamiltonian systems.
Arch. Rat. Mech. Anal. {\bf 78} (1982) 315--333

\bibitem {clar:eke:2}
Clarke, F., Ekeland, I.:
Solutions p\'{e}riodiques, du
p\'{e}riode donn\'{e}e, des \'{e}quations hamiltoniennes.
Note CRAS Paris {\bf 287} (1978) 1013--1015

\bibitem {mich:tar}
Michalek, R., Tarantello, G.:
Subharmonic solutions with prescribed minimal
period for nonautonomous Hamiltonian systems.
J. Diff. Eq. {\bf 72} (1988) 28--55

\bibitem {tar}
Tarantello, G.:
Subharmonic solutions for Hamiltonian
systems via a $\bbbz_{p}$ pseudoindex theory.
Annali di Matematica Pura (to appear)

\bibitem {rab}
Rabinowitz, P.:
On subharmonic solutions of a Hamiltonian system.
Comm. Pure Appl. Math. {\bf 33} (1980) 609--633

% JasonDWilliams,“Web-stylerankingandslucombina- tion for dialog state tracking,” in 15th Annual Meeting of the Special Interest Group on Discourse and Dialogue, 2014, p. 282.

\end{thebibliography}
\end{document}
